\section{Conclusions}

Across 472,530 named-basin GLODAP v2.2023 observations, the empirical dependence of $\tco$ on salinity, temperature, and apparent oxygen utilisation was substantially compressible into a five-parameter rank-1 quadratic projection. The model improved on a linear baseline in every basin, and adding a second squared projection produced no consistent validation benefit. More complex symbolic expressions achieved lower validation error, but at two- to three-fold greater expression complexity.

The compact structure was not invariant across symbolic operator spaces or validation regimes. Performance also depended on basin, depth, and water-mass context: AAIW showed large positive bias, and a globally fitted Southern Ocean model failed in the abyss while a local refit performed well. The appropriate conclusion is therefore an empirical and conditional one. Rank-1 quadratic structure is a useful low-complexity approximation for the sampled ocean-carbon relationship, but it is neither a universal predictor nor evidence that the physical dimensionality of ocean carbon dynamics is one.
